% strat_t1_performance — OOS performance of the GCF bond-timing strategy.
% Numbers transcribed from tables/strat_t1_performance.pdf (generated by strategy.R).
\small
\begin{tabular}{l c c c c c}
  \toprule
  Strategy & Mean (\%) & Vol (\%) & Sharpe & CER, $\gamma=5$ (\%) & CER, $\gamma=10$ (\%) \\
  \midrule
  $\GCF$-timed           & $3.02$ & $7.89$ & $0.38$ & $1.46$ & $-0.10$ \\
  Recursive-mean timing  & $2.39$ & $7.00$ & $0.34$ & $1.16$ & $-0.07$ \\
  Buy-and-hold           & $2.17$ & $6.98$ & $0.31$ & $0.95$ & $-0.27$ \\
  \bottomrule
\end{tabular}
\tabnotes{This table reports the real-time performance of timing a 10Y
  GDP-weighted global bond portfolio with the global cycle factor over
  2005--2024 ($n=221$ months; the portfolio's out-of-sample $R^{2}$ against the
  recursive mean is $0.125$). Exposure follows the mean--variance rule
  \eqref{eq:strat-weight},
  $w_{t}\propto\widehat{\E}_{t}[\rx_{t+12}]/\widehat{\sigma}^{2}_{t}$, with the
  forecast taken from the recursive $\GCF^{\mathrm{oos}}$ regression, so the
  strategy is doubly out-of-sample. The Sharpe ratio is scale-invariant; mean,
  volatility, and the certainty-equivalent (CER) return are reported at equal
  average exposure across strategies so that the variance penalty is
  comparable. CER is the annual certainty equivalent of a mean--variance
  investor with risk aversion $\gamma$. Returns are local-currency
  (currency-hedged) excess returns based on overlapping twelve-month holding
  periods.}
